\exercisegroup

\begin{enumerate}
\def\labelenumi{\arabic{enumi})}
\tightlist
\item
  Let E and F be two infinite dimensional Hilbert spaces satisfying the first axiom of countability, \((a_{n})\) an orthonormal basis of E, and \((b_{n})\) an orthonormal basis of F.
\end{enumerate}

\begin{enumerate}
\def\labelenumi{\alph{enumi})}
\tightlist
\item
  Let u be a continuous linear mapping from E into F; let \(u(a_{n}) = \sum_{m} \alpha_{mn} b_{m}\) . Show that
\[
\sum_ {n} | \alpha_ {m n} | ^ {2} \leqslant \| u \| ^ {2}
\]

\[
\sum_ {m} | \alpha_ {m n} | ^ {2} \leqslant \| u \| ^ {2}
\]
for every \(m\) and \(n\) .

\begin{enumerate}
\def\labelenumi{\alph{enumi})}
\setcounter{enumi}{1}
\tightlist
\item
  Give an example of a double sequence \((\alpha_{mn})\) such that \(\sum_{m} |\alpha_{mn}|^{2} \leqslant 1\) for all \(n\) and \(\sum_{n} |\alpha_{mn}|^{2} \leqslant 1\) for all \(m\) , but such that there does not exist any continuous linear mapping \(u\) from E into F such that \(\langle u(a_n)|b_m\rangle = \alpha_{mn}\) for every pair of integers \((m, n)\) . (Show that if I \(\subset\) N is a set of \(p\) integers, and if V \(_p\) (resp. W \(_p\) ) is the subspace of E (resp. F) generated by the \(a_n\) (resp. \(b_n\) ) such that \(n \in I\) , then there exists a linear mapping \(u_p\) from V \(_p\) onto W \(_p\) such that \(\langle u_p(a_n)|b_m\rangle = \frac{1}{\sqrt{p}}\) for \(m \in I\) and \(n \in I\) , and that \(\|u_p\| \geqslant \sqrt{p}\) .)
\end{enumerate}

\begin{enumerate}
\def\labelenumi{\arabic{enumi})}
\setcounter{enumi}{1}
\tightlist
\item
  Let \(A = (\alpha_{mn})_{(m,n)\in \mathbf{N}\times \mathbf{N}}\) be a double sequence of complex numbers, which we also call an infinite matrix. For every point \(x = (x_{n})\) of the direct sum space \(\mathbf{C}^{(\mathbf{N})}\) , the sums \(y_{m} = \sum_{n}\alpha_{mn}x_{n}\) are defined; let \(A.x\) be the point \((y_{m})\) of the product space \(\mathbf{C}^{\mathbf{N}}\) , then \(x\mapsto A.x\) is a linear mapping from \(\mathbf{C}^{(\mathbf{N})}\) into \(\mathbf{C}^{\mathbf{N}}\) , and every linear mapping from \(\mathbf{C}^{(\mathbf{N})}\) into \(\mathbf{C}^{\mathbf{N}}\) is of this form. Let \(E_{n}\) denote the subspace of \(\mathbf{C}^{(\mathbf{N})}\) generated by the first \(n\) vectors of the canonical basis, \(P_{n}\) the canonical projection from \(\mathbf{C}^{\mathbf{N}}\) onto \(E_{n}\) ; when \(E_{n}\) is assigned the norm induced by that of the space \(\ell_{\mathbf{C}}^{2}(\mathbf{N})\) , \(\| u\|\) denotes the norm of the linear mapping \(u\) from the finite dimensional hilbertian space \(E_{n}\) into itself.
\end{enumerate}

\begin{enumerate}
\def\labelenumi{\alph{enumi})}
\item
  In order that the image of \(\mathbf{C}^{(\mathbf{N})}\) under the mapping \(x\mapsto A.x\) be contained in \(\ell_{\mathbf{C}}^2 (\mathbf{N})\) and that this mapping extend to a continuous linear mapping from \(\ell_{\mathbf{C}}^2 (\mathbf{N})\) into itself, it is necessary and sufficient that the norms \(\| P_nAP_n\|\) are bounded. This implies that the rows and the columns of \(A\) belong to \(\ell_{\mathbf{C}}^2 (\mathbf{N})\) (exerc. 1).
\item
  Let \(A^*\) denote the infinite matrix \((\alpha_{mn}')\) , where \(\alpha_{mn}' = \overline{\alpha}_{nm}\) . If the columns of \(A\) belong to \(\ell_{\mathbf{C}}^2(\mathbf{N})\) (in other words, if \(x \mapsto A.x\) maps \(\mathbf{C}^{(\mathbf{N})}\) into \(\ell_{\mathbf{C}}^2(\mathbf{N})\) ), then the series \(\beta_{mn} = \sum_{p} \overline{\alpha}_{pm} \alpha_{pn}\) are
\end{enumerate}

absolutely convergent, and we put \(A^{*}A = (\beta_{mn})\) . Show that for \(x \mapsto A.x\) to extend to a continuous linear mapping \(u\) from \(\ell_{\mathbf{C}}^{2}(\mathbf{N})\) into itself, it is necessary and sufficient that the norms \(\| P_n(A^*A)P_n\|\) are bounded (we have \(\langle P_n(A^*A)P_n.x|x\rangle = \| AP_n .x\|^2\) for all \(x\in E_n\) ). Then \(x\mapsto A^{*}A.x\) extends to a positive hermitian mapping \(u^{*}u\) from \(\ell_{\mathbf{C}}^{2}(\mathbf{N})\) into itself.

\begin{enumerate}
\def\labelenumi{\alph{enumi})}
\setcounter{enumi}{2}
\item
  For two infinite matrices \(X = (\xi_{mn})\) , \(Y = (\eta_{mn})\) , we say that the product \(XY\) is defined if the series \(\zeta_{mn} = \sum_{p} \xi_{mp} \eta_{pn}\) are absolutely convergent, and then we put \(XY = (\zeta_{mn})\) . We say that a power \(X^k\) ( \(k\) integer \(> 1\) ) is defined if \(X^{k-1}\) and \(X^{k-1}X\) are defined and then we put \(X^k = X^{k-1}X\) ; in this case, \(X^pX^q = X^k\) for every pair of integers \(p, q\) such that \(p + q = k\) . If \(A\) is an infinite matrix whose columns are in \(\ell_{\mathbf{C}}^2(\mathbf{N})\) and if the product \((A^*A)^2\) is defined, show that for all \(x \in E_n\) , we have \(\langle (P_nA^*AP_n)^2.x|x\rangle \leqslant \langle P_n(A^*A)^2P_n.x|x\rangle\) , and deduce that \(\| P_nA^*AP_n\|^2 \leqslant \| P_n(A^*A)^2P_n\|\) .
\item
  In order that an infinite matrix A be such that the image of \(\mathbf{C}^{(\mathbf{N})}\) under \(x\mapsto A.x\) is contained in \(\ell_{\mathbf{C}}^2 (\mathbf{N})\) and that \(x\mapsto A.x\) extends to a continuous linear mapping from \(\ell_{\mathbf{C}}^2 (\mathbf{N})\) into itself, it is necessary and sufficient that the following three conditions are satisfied:
\end{enumerate}

\begin{enumerate}
\def\labelenumi{(\roman{enumi})}
\item
  the rows and columns of \(A\) are in \(\ell_{\mathbf{C}}^{2}(\mathbf{N})\) ;
\item
  the powers \((A^{*}A)^{k}\) are defined for every integer \(k > 1\) ;
\item
  we have
\end{enumerate}

\[
\sup _ {n} \left(\sup _ {m} | ((A ^ {*} A) _ {m m} ^ {n}) ^ {1 / n} |\right) <   + \infty
\]

where \((A^{*}A)_{mm}^{n}\) denotes the term with indices \((m,m)\) of the matrix \((A^{*}A)^{n}\) . (Observe that if \(C\) is the matrix, with respect to the canonical basis of \(\mathrm{E}_n\) , of a positive hermitian endomorphism of \(\mathrm{E}_n\) , then \(\| C\| \leqslant n\sup_{1\leqslant i\leqslant n}|C_{ii}|\) , by considering the trace of \(C\) and by diagonalizing \(C\) . Using the inequality proved in \(c\) ), show that

\[
\left\| P _ {n} A ^ {*} A P _ {n} \right\| \leqslant n ^ {2 - k} \sup _ {1 \leqslant i \leqslant n} \left| ((A ^ {*} A) _ {i i} ^ {2 k}) \right| ^ {2 - k}
\]

for every integer \(k > 1\) , if conditions (i), (ii) and (iii) are satisfied.)

¶ 3) Let \((a_{ij})_{(i,j)\in\mathbb{I}\times\mathbb{I}}\) be a countably infinite double family of complex numbers. Assume that there exist two numbers \(\beta > 0\) , \(\gamma > 0\) , and a family \((p_i)_{i\in\mathbb{I}}\) of numbers > 0 satisfying the relations

\[
\sum_ {i} p _ {i} | a _ {i j} | \leqslant \beta p _ {j}, \quad \sum_ {j} | a _ {i j} | p _ {j} \leqslant \gamma p _ {i}\tag{*}
\]

for all \(i, j\) in \(\mathbb{I}\) .
a) Show that there exists a continuous endomorphism \(u\) from \(\ell_{\mathbf{C}}^2 (\mathrm{I})\) , with norm \(\leqslant (\beta \gamma)^{1 / 2}\) , such that for all \(x = (x_{i})_{i\in \mathbf{I}}\) in \(\ell_{\mathbf{C}}^2 (\mathrm{I})\) , we have \(u(x) = y\) , where \(y = (y_{i})\) is given by \(y_{i} = \sum_{j}a_{ij}x_{j}\) (for \(x = (x_{i})\) and \(y = (y_{i})\) in \(\mathbf{C}^{(\mathbb{I})}\) , put \(v_{ij} = |x_i| (p_j|a_{ij}| / p_i)^{1 / 2}\) , \(w_{ij} = |y_j|(p_i|a_{ij}| / p_j)^{1 / 2}\) , and find a bound for \(\sum_{i,j}v_{ij}w_{ij}\) ).

\begin{enumerate}
\def\labelenumi{\alph{enumi})}
\setcounter{enumi}{1}
\tightlist
\item
  For I, take the set of all integers \(\geqslant 1\) , and let \(a_{ij} = (i + j)^{-1}\) . Show that the conditions (*) are satisfied with \(p_i = i^{-1/2}\) and \(\beta = \gamma = \pi\) (these being the best possible constants) (compare the series in (*) with an integral). In an analogous manner, treat the case where I = N and \(a_{ij} = (i + j + 1)^{-1}\) (« Hilbert's matrix »).
\end{enumerate}

* c) Let \(\mathcal{H} = \mathrm{H}^2 (\mathbf{D})\) be the Hardy space, consisting of all functions \(f(z) = \sum_{n=0}^{\infty} a_n z^n\) , analytic in the open disc \(\mathbf{D}:|z| < 1\) and such that \(\| f\|^2 = \sum_n |a_n|^2 < +\infty\) ; then \(\| f\|\) is a norm on \(\mathcal{H}\) , for which \(\mathcal{H}\) is isomorphic to \(\ell_{\mathbf{C}}^2 (\mathbf{N})\) . Given two functions \(f, g\) in \(\mathcal{H}\) , show that the function \(t \mapsto f(t) g(t)\) of the real variable \(t\) is integrable on \([0,1]\) with respect to the Lebesgue measure and that the formula \(\mathbf{B}(f,g) = \int_0^1 f(t) g(t) dt\) defines a continuous bilinear form on \(\mathcal{H} \times \mathcal{H}\) (consider \(\mathbf{B}(f,f) = \int_0^1 f(t)^2 dt\) for a function \(f \in \mathcal{H}\) of the form \(f(z) = \sum_{n=0}^{\mathbf{N}} a_n z^n\) with \(a_n \geqslant 0\) for \(0 \leqslant n \leqslant \mathbf{N}\) ; use Cauchy's theorem to establish the relation

\[
\int_ {- 1} ^ {1} f (t) ^ {2} d t = - i \int_ {0} ^ {\pi} f (e ^ {i \theta}) ^ {2} e ^ {i \theta} d \theta
\]

from which we get \(\mathbf{B}(f, f) \leqslant \frac{1}{2} \int_{-\pi}^{\pi} |f(e^{i\theta})|^2 d\theta = \pi \| f\|^2\) . In this way, get the result of \(b\) , according to which the Hilbert's matrix defines an endomorphism of norm \(\leqslant \pi\) of the hilbertian space \(\ell_{\mathbf{C}}^2(\mathbf{N})\) . *

\begin{enumerate}
\def\labelenumi{\arabic{enumi})}
\setcounter{enumi}{3}
\tightlist
\item
  Let \(\mathbf{E}\) be a complex hilbertian space of finite dimension \(d\) .
\end{enumerate}

\begin{enumerate}
\def\labelenumi{\alph{enumi})}
\item
  For every \(u \geqslant 0\) in \(\mathcal{L}(\mathrm{E})\) (V, p.~45), prove that there exists a unique \(v \geqslant 0\) in \(\mathcal{L}(\mathrm{E})\) such that \(u = v^2\) (diagonalize \(u\) ); we write \(v = u^{1/2}\) .
\item
  For every \(u \in \mathcal{L}(\mathrm{E})\) , put \(\mathrm{abs}(u) = (u^{*}u)^{1/2}\) . Show that u and \(\mathrm{abs}(u)\) have the same norm and that we have \(\mathrm{abs}(\symbf{\Lambda}^{n}(u)) = \symbf{\Lambda}^{n}(\mathrm{abs}(u))\) for every integer \(n \leqslant d\) .
\item
  Let \(s_{1}(u) \geqslant s_{2}(u) \geqslant \ldots \geqslant s_{d}(u) \geqslant 0\) be the sequence of eigenvalues of \(\mathrm{abs}(u)\) counted with their order of multiplicity. Show that \(\|u\| = s_{1}(u)\) and that for every integer \(n \leqslant d\) , \(\| \symbf{\Lambda}^{n}(u)\| = s_{1}(u) s_{2}(u) \ldots s_{n}(u)\) . In order that \(\| \symbf{\Lambda}^{n}(u)\| = \| u\|^{n}\) for all n such that \(1 \leqslant n \leqslant d\) , it is necessary and sufficient that \(u^{*}u\) is a homothety, in other words, that u is a scalar multiple of a unitary operator.
\end{enumerate}

\begin{enumerate}
\def\labelenumi{\arabic{enumi})}
\setcounter{enumi}{4}
\tightlist
\item
  Let E, F be two Hilbert spaces, and u a continuous linear mapping from E into F. Let \(\ell(u)\) denote the set of all \(x \in E\) such that \(\|u(x)\| = \|u\|. \|x\|\) .
\end{enumerate}

\begin{enumerate}
\def\labelenumi{\alph{enumi})}
\item
  Show that \(\ell(u)\) is the closed vector subspace of E which is the kernel of \(u^{*}u - \| u\|^{2}1_{\mathrm{E}}\) and is orthogonal to the kernel of \(u\) .
\item
  Show that the restriction of \(u\) to \(\ell(u)\) is a bijection from \(\ell(u)\) onto \(\ell(u^*)\) , whose inverse bijection is the restriction of \(\|u\|^{-2}.u^*\) to \(\ell(u^*)\) ; moreover the image of the orthogonal complement \((\ell(u))^{\circ}\) under \(u\) is contained in \((\ell(u^*))^{\circ}\) . If \(u_1\) is the restriction of \(u\) to \((\ell(u))^{\circ}\) , considered as a mapping from \((\ell(u))^{\circ}\) into \((\ell(u^*))^{\circ}\) , the adjoint \(u_1^*\) is the restriction of \(u^*\) to \((\ell(u^*))^{\circ}\) ; if \(\ell(u) \neq E\) , let \(\langle u\rangle\) , called the sub-norm of \(u\) , denote the norm \(\|u_1\|\) ; if \(\ell(u) = E\) , we put \(\langle u\rangle = 0\) . Then \(\langle u^*\rangle = \langle u\rangle\) .
\end{enumerate}

\begin{enumerate}
\def\labelenumi{\arabic{enumi})}
\setcounter{enumi}{5}
\tightlist
\item
  Let E be a hilbertian space, M, N two closed subspaces of E and \(M^{\circ}\) , \(N^{\circ}\) their respective orthogonal complements; let \(p_{M}\) , \(p_{N}\) denote the orthogonal projections on M and N respectively, such that \(1_{E} - p_{M}\) , \(1_{E} - p_{N}\) are the orthogonal projections on \(M^{\circ}\) and \(N^{\circ}\) respectively. Put \(u_{\mathrm{NM}} = (1_{\mathrm{E}} - p_{\mathrm{N}}) p_{\mathrm{M}}\) and \(\delta(\mathbf{M}, \mathbf{N}) = \|u_{\mathrm{NM}}\|\) ; we have \(\delta(\mathbf{M}, \mathbf{N}) = \delta(\mathbf{N}^{\circ}, \mathbf{M}^{\circ}) \leqslant 1\) ; the relation \(\delta(\mathbf{M}, \mathbf{N}) < 1\) implies \(M \cap N^{\circ} = \{0\}\) .
\end{enumerate}

\begin{enumerate}
\def\labelenumi{\alph{enumi})}
\item
  Let \(\tilde{\mathbf{M}}\) denote the orthogonal complement in \(\mathbf{M}\) of \(\mathbf{M} \cap \mathbf{N}^{\circ}\) , and let \(\varepsilon(\mathbf{M}, \mathbf{N}) = \delta(\tilde{\mathbf{M}}, \mathbf{N})\) . Show that (with the notations of exerc. 5) \(\ell(u_{\mathrm{NM}}) = \mathbf{M} \cap \mathbf{N}^{\circ}\) and deduce that \(\varepsilon(\mathbf{M}, \mathbf{N}) = \langle u_{\mathrm{NM}} \rangle \leqslant \delta(\mathbf{M}, \mathbf{N})\) ; moreover, if \(\mathbf{M} \cap \mathbf{N}^{\circ} = \{0\}\) (and in particular if \(\delta(\mathbf{M}, \mathbf{N}) < 1\) ), then \(\varepsilon(\mathbf{M}, \mathbf{N}) = \delta(\mathbf{M}, \mathbf{N})\) .
\item
  For a continuous linear mapping \(u\) from E into itself, we designate by conorm of \(u\) the number \(c(u) = \inf \|u(x)\|/\|x\|\) , where \(x\) ranges over the set of all vectors \(\neq 0\) orthogonal to \(u^{-1}(0)\) (if \(u = 0\) , put \(c(u) = 1\) ). For \(u(\mathrm{E})\) to be closed in E, it is necessary and sufficient that \(c(u) > 0\) (I, p.~17, th. 1). We have \(c(u^{*}) = c(u)\) .
\item
  Let \(v_{\mathbf{NM}} = p_{\mathbf{N}}p_{\mathbf{M}}\) . Show that
\end{enumerate}

\[
\varepsilon (\mathbf {M}, \mathbf {N}) ^ {2} + c (v _ {\mathbf {N M}}) ^ {2} = 1
\]

(observe that \(\|u_{\mathrm{NM}}(x)\|^{2} + \|v_{\mathrm{NM}}(x)\|^{2} = \|p_{\mathrm{M}} \cdot x\|^{2}\) , and that the kernel of \(v_{NM}\) is \(M^{\circ} + (M \cap N^{\circ})\) and deduce that \(\langle u_{NM}\rangle^{2} \leqslant 1 - c(v_{\mathrm{NM}})^{2}\) .)
\item
  Deduce from b) and c) that \(\varepsilon(\mathrm{N}, \mathrm{M}) = \varepsilon(\mathrm{M}, \mathrm{N})\) and, using a), that \(\varepsilon(\mathrm{M}^{\circ}, \mathrm{N}^{\circ}) = \varepsilon(\mathrm{M}, \mathrm{N})\) .
\item
  Put \(g(\mathbf{M}, \mathbf{N}) = \|p_{\mathbf{M}} - p_{\mathbf{N}}\|\) (cf.~V, p.~64, exerc. 17). Show that

\[
g (\mathbf {M}, \mathbf {N}) = \sup (\delta (\mathbf {M}, \mathbf {N}), \delta (\mathbf {N}, \mathbf {M}))
\]

(observe that \(p_{\mathbf{M}} - p_{\mathbf{N}} = (1_{\mathrm{E}} - p_{\mathbf{N}})p_{\mathbf{M}} - p_{\mathbf{N}}(1_{\mathrm{E}} - p_{\mathbf{M}})\) ); deduce that \(\varepsilon (\mathbf{M},\mathbf{N})\leqslant g(\mathbf{M},\mathbf{N})\) . If \(\mathbf{M}\cap \mathbf{N}^{\circ} = \mathbf{N}\cap \mathbf{M}^{\circ} = \{0\}\) , we have the relation \(\varepsilon (\mathbf{M},\mathbf{N}) = \delta (\mathbf{M},\mathbf{N}) = \delta (\mathbf{N},\mathbf{M}) = g(\mathbf{M},\mathbf{N})\) . If \(g(\mathbf{M},\mathbf{N}) < 1\) , then \(\mathbf{M}\cap \mathbf{N}^{\circ} = \mathbf{N}\cap \mathbf{M}^{\circ} = \{0\}\) .

\end{enumerate}

\begin{enumerate}
\def\labelenumi{\alph{enumi})}
\setcounter{enumi}{5}
\tightlist
\item
  Let \(Q_{M}\) , \(Q_{N}\) be two continuous projections in E, with respective images M and N; give another proof of the relation \(g(\mathrm{M}, \mathrm{N}) \leqslant \|Q_{\mathrm{M}} - Q_{\mathrm{N}}\|\) (V, p.~64, exerc. 17). (Note that for all \(x \in E\) , we have \(\|(1_{E} - Q_{M}).x\|^{2} + \|Q_{M}^{*}.x\|^{2} = \|x\|^{2} + \|(Q_{\mathrm{M}} - Q_{\mathrm{M}}^{*}).x\|^{2}\) , and apply this relation to \(x = (p_{\mathrm{M}} - p_{\mathrm{N}}).y\) , observing the relations \((1_{\mathrm{E}} - Q_{\mathrm{M}})(p_{\mathrm{M}} - p_{\mathrm{N}}) = (Q_{\mathrm{M}} - Q_{\mathrm{N}}) p_{\mathrm{N}}\) and \((p_{\mathrm{M}} - p_{\mathrm{N}}) Q_{\mathrm{M}} = (1_{\mathrm{E}} - p_{\mathrm{N}})(Q_{\mathrm{M}} - Q_{\mathrm{N}}).\) )
\end{enumerate}

¶ 7) a) With the notations of exerc. 6, show that, for \(\mathbf{M} + \mathbf{N}^{\circ}\) to be closed, it is necessary and sufficient that \(\mathbf{M} + \mathbf{N}^{\circ} = (\mathbf{M}^{\circ} \cap \mathbf{N})^{\circ}\) .

\begin{enumerate}
\def\labelenumi{\alph{enumi})}
\setcounter{enumi}{1}
\tightlist
\item
  Show that the following properties are equivalent :
\end{enumerate}

\(\alpha)\)

\[
\varepsilon (\mathbf {M}, \mathbf {N}) <   1;
\]

β) \(M + N^{\circ}\) is closed in E;

\(\gamma)\) If \(\tilde{\mathbf{M}}\) is the orthogonal complement of \(\mathbf{M} \cap \mathbf{N}^{\circ}\) in \(\mathbf{M}\) and \((\mathbf{N}^{\circ})^{\sim}\) that of \(\mathbf{M} \cap \mathbf{N}^{\circ}\) in \(\mathbf{N}^{\circ}\) , then \(\mathbf{E}\) is the direct sum of \(\mathbf{M}^{\circ} \cap \mathbf{N}\) , \(\mathbf{M} \cap \mathbf{N}^{\circ}\) , \(\tilde{\mathbf{M}}\) and \((\mathbf{N}^{\circ})^{\sim}\) . Moreover, if \(R\) and \(S\) are respectively the projections from \(\mathbf{E}\) onto \(\mathbf{M}\) and \((\mathbf{N}^{\circ})^{\sim}\) corresponding to this decomposition, then \(\| R \| = \| S \| = (1 - \varepsilon^2 (\mathbf{M}, \mathbf{N}))^{1/2}\) . (To see that \(\alpha\) ) implies \(\beta\) ) observe first that if \(x \in \tilde{\mathbf{M}}\) and \(y \in (\mathrm{N}^\circ)^{\sim}\) , then \(|\langle x | y \rangle| \leqslant \varepsilon (\mathbf{M}, \mathbf{N}) \| x \| . \| y \|\) ; then, let \(u = x + y + t\) be the decomposition of an element \(u \in \mathbf{M} + \mathrm{N}^\circ\) with \(x \in \mathbf{M}\) , \(y \in (\mathrm{N}^\circ)^{\sim}\) and \(t \in \mathbf{M} \cap \mathrm{N}^\circ\) , deduce that \(\| x \| \leqslant (1 - \varepsilon^2 (\mathbf{M}, \mathbf{N}))^{-1/2} \| u \|, \| y \| \leqslant (1 - \varepsilon^2 (\mathbf{M}, \mathbf{N}))^{-1/2} \| u \|\) and \(\| t \| \leqslant \| u \|\) . To prove that \(\gamma\) ) implies \(\alpha\) ), consider the decomposition \(v = v_1 + v_2\) for a \(v \in \tilde{\mathbf{M}}\) , where \(v_1\) is the orthogonal projection of \(v\) onto \(\tilde{\mathbf{N}}\) , the orthogonal complement of \(\mathbf{N} \cap \mathbf{M}^\circ\) in \(\mathbf{N}\) ; we have \(R.v_1 = v\) , and if we had \(\varepsilon (\mathbf{M}, \mathbf{N}) = 1\) , there would exist a sequence \((v_n) \in \mathbf{M}\) such that \(\| v_n \| = 1\) and such that \(\| (v_n)_1\|\) tends to 0. Next show that the restriction \(R_1\) of \(R\) to \(\tilde{\mathbf{N}}\) is a bijection from \(\tilde{\mathbf{N}}\) onto \(\tilde{\mathbf{M}}\) ; to calculate \(\| R \|\) , show that \(\| R_1^{-1} \| \leqslant (1 - \varepsilon^2 (\mathbf{M}, \mathbf{N}))^{1/2}\) .

\begin{enumerate}
\def\labelenumi{\alph{enumi})}
\setcounter{enumi}{2}
\tightlist
\item
  Deduce from b) that if \(M + N^{\circ}\) is closed, then so is \(M^{\circ} + N\) .
\end{enumerate}

\begin{enumerate}
\def\labelenumi{\arabic{enumi})}
\setcounter{enumi}{7}
\tightlist
\item
  \begin{enumerate}
  \def\labelenumii{\alph{enumii})}
  \tightlist
  \item
    Let E be a hilbertian space, and let T be a continuous linear mapping from E into itself such that \(\|T\|\leqslant1\) . Show that the relations \(T.x=x,\langle T.x|x\rangle=\|x\|\) , \(T^{*}.x=x\) are equivalent, and that the kernel of \(1_{E}-T\) and the closure of the image of \(1_{E}-T\) are the orthogonal complements.
  \end{enumerate}
\end{enumerate}

\begin{enumerate}
\def\labelenumi{\alph{enumi})}
\setcounter{enumi}{1}
\tightlist
\item
  Let \(T\) be a continuous linear mapping from \(E\) into itself, satisfying the inequality

\[
\| x - T. x \| ^ {2} \leqslant \| x \| ^ {2} - \| T. x \| ^ {2}\tag{1}
\]

for all \(x \in E\) . Then \(\|T.x\| < \|x\|\) for all x such that \(T.x \neq x\) ; for every \(x \in E\) , the sequence \((T^{n}.x)\) converges to a point P.x, and P is the orthogonal projector onto the kernel of \(1_{E} - T\) .
\item
  Let \(P_{1}, \ldots, P_{r}\) be orthoprojectors on E. Show that the product \(T = P_{1}P_{2} \ldots P_{r}\) satisfies relation (1) (argue by induction on r); hence the orthoprojector P is the orthogonal projector onto the intersection of the images of the projectors \(P_{j}\) (note that if \(\|P_{j}.x\| < \|x\|\) for an index j, then \(\|T.x\| < \|x\|\) ).
\end{enumerate}

\begin{enumerate}
\def\labelenumi{\arabic{enumi})}
\setcounter{enumi}{8}
\tightlist
\item
  Let E be a hilbertian space and \((P_{j})_{j\in\mathbb{N}}\) be a sequence of orthoprojectors an E, such that, for all \(j\in N\) , there exists an \(n_{j}\in N\) such that for all \(k\in N\) , at least one of the orthoprojectors \(P_{k}, P_{k+1}, ..., P_{k+n_{j}}\) is equal to \(P_{j}\) . Put \(R_{s}=P_{s}P_{s-1}\ldots P_{0}\) for all \(s\in N\) .
\end{enumerate}

\begin{enumerate}
\def\labelenumi{\alph{enumi})}
\tightlist
\item
  For every \(x\in E\) , let \(x_{s}=R_{s}\cdot x\) . Show that \(\sum_{s}\|x_{s-1}-x_{s}\|^{2}\leqslant\|x\|^{2}\) , and deduce that for
\end{enumerate}

every integer \(r \geqslant 1\) , \(x_{s + r} - x_s\) tends to 0 as \(s\) tends to \(+\infty\) .

\begin{enumerate}
\def\labelenumi{\alph{enumi})}
\setcounter{enumi}{1}
\item
  Let \((x_{s_k})\) be a sequence extracted from \((x_s)\) which tends to a limit \(y\) weakly; then every sequence \((x_{s_k + r})\) also tends to \(y\) weakly. Deduce that \(y\) belongs to each of the subspaces \(\mathbf{M}_j = P_j(\mathbf{E})\) . (For each \(j\) , there exists \(r_k\) such that \(0 \leqslant r_k \leqslant n_j\) and \(s_k + r_k = j\) ; show that the sequence \((x_{s_k + r_k})\) tends to \(y\) weakly.)
\item
  Show that the sequence \((x_{s})\) converges weakly to the orthogonal projection from \(x\) onto the intersection \(\mathbf{M}\) of the \(\mathbf{M}_{j}\) . (Reduce to the case where \(\mathbf{M} = \{0\}\) , and use \(b\) ) and the weak compactness of every closed ball in E.)
\end{enumerate}

¶ 10) Let E be a hilbertian space, u a positive endomorphism of E.

\begin{enumerate}
\def\labelenumi{\alph{enumi})}
\tightlist
\item
  Show that for all \(x \in E\) , we have

\[
\left\| u (x) \right\| ^ {2} \leqslant \| u \|. \langle u (x) | x \rangle
\]

(observe that \(\langle u(x)|u(x)\rangle^2\leqslant \langle u(x)|x\rangle \langle u^2 (x)|u(x)\rangle\) , by V, p.~3, prop. 2).
\item
  Let M be a closed vector subspace of E, \(\mathbf{M}^{\circ}\) its orthogonal complement. Let \(x\in \mathbf{M}\) , and let \(f(x)\) be the lower bound of \(\langle u(x + y)|x + y\rangle\) as \(y\) ranges over \(\mathbf{M}^{\circ}\) . For every \(\varepsilon >0\) , let \(\operatorname {E}(x,\varepsilon)\) be the set of all \(y\in \mathbf{M}^{\circ}\) such that \(\langle u(x + y)|x + y\rangle \leqslant f(x) + \varepsilon\) . Show that \(\operatorname {E}(x,\varepsilon)\) is convex and that for all \(z\in \mathbf{M}^{\circ}\) , we have

\[
\langle u (x + y) | z \rangle^ {2} \leqslant \varepsilon \langle u (z) | z \rangle \quad \text { for } \quad y \in \mathrm{E} (x, \varepsilon)\tag{*}
\]

(consider the function \(g:t \mapsto \langle u(x + y + tz)/x + y + tz \rangle\) of the real variable t, which attains its minimum at a point \(t_{0}\) and note that \(g(t_{0}) \geqslant f(x)\) and \(g(0) \leqslant f(x) + \varepsilon\) ).
\item
  For every integer \(n \geqslant 1\) , let \(y_{n} \in \mathrm{E}(x, 1/n)\) ; show that the sequence \((u(x + y_{n}))\) tends to a limit \(x_{1}\) belonging to M, and that the sequence \((\langle u(y_{n})|y_{n}\rangle)\) is bounded (find a bound for the numbers \(\langle u(y_{n} - y_{m})|y_{n} - y_{m}\rangle\) for \(m \geqslant n\) and \(|\langle u(x + y_{n})|z\rangle|\) for \(z \in M^{\circ}\) using inequality (*)).

\item
  Let \((y_n')\) be a sequence of points of \(\mathbf{M}^\circ\) such that \(\langle u(y_n' - y_m')|y_n' - y_m'\rangle\) is arbitrarily small when \(m\) and \(n\) are large enough, and such that the sequence \((u(x + y_n'))\) has a limit \(x_1' \in \mathbf{M}\) ;

show that \(x_1' = x_1\) . (Let \(Q(z) = \langle u(z)|z \rangle\) ; first show that the number \(Q((y_p - y_p') - (y_q - y_q'))\) is arbitrarily small as soon as \(p\) and \(q\) are large enough, and deduce that the sequence \((Q(y_n - y_n'))\) is bounded; using the fact that \(\langle x_1' - x_1|y_p - y_p' \rangle = 0\) for all \(p\) , show that the sequence \((Q(y_n - y_n'))\) tends to 0 and use \(a\) ).)

\item
  Deduce from \(d\) ) that the point \(x_{1}\) does not depend on the choice of the \(y_{n} \in \mathrm{E}(x, 1/n)\) , and that if we put \(u_{1}(x) = x_{1}\) , then \(u_{1}\) is a linear mapping from \(\mathbf{M}\) into itself. Show that \(0 \leqslant \langle u_{1}(x)|x\rangle \leqslant \langle u(x)|x\rangle\) for all \(x \in \mathbf{M}\) and consequently that \(u_{1}\) is continuous and is an endomorphism of \(\mathbf{M}\) which is \(\geqslant 0\) . (Observe that \(\langle u(x + y_n)|y_n\rangle\) tends to 0 and \(\langle u(x + y_n)|x + y_n\rangle\) tends to \(f(x)\) .)
\item
  Let \(p_{M}\) be the orthoprojector with image M, and let \(u_{0} = u_{1} \circ p_{M}\) . We have \(0 \leqslant u_{0} \leqslant u\) , and M is stable under \(u_{0}\) , and the restriction of \(u_{0}\) to \(M^{\circ}\) is null. Show that \(u_{0}\) is the largest element in the family of all endomorphisms \(v \geqslant 0\) such that \(v \leqslant u\) , that M is stable under v and is such that the restriction of v to \(M^{\circ}\) is null.
\end{enumerate}

\begin{enumerate}
\def\labelenumi{\arabic{enumi})}
\setcounter{enumi}{10}
\tightlist
\item
  Let \(\mathrm{E}\) and \(\mathrm{F}\) be two hilbertian spaces. Show that for every element \(u\) in \(\mathrm{E} \hat{\otimes}_2 \mathrm{F}\) , there exists an orthonormal sequence \((e_n)\) in \(\mathrm{E}\) , an orthonormal sequence \((f_n)\) in \(\mathrm{F}\) and a sequence \((\lambda_n)\) of numbers \(\geqslant 0\) such that \(\sum_{n} \lambda_n^2 < +\infty\) and that \(u = \sum_{n} \lambda_n e_n \otimes f_n\) ; then \(\| u \|_2^2 = \sum_{n} \lambda_n^2\)
\end{enumerate}

(cf.~V, p.~55, th. 2 and p.~53, th. 1).

¶ 12) Let E be a real hilbertian space and V be a closed convex cone in E, with vertex 0, let V° be the polar cone of V (in E, identified canonically with its dual).

\begin{enumerate}
\def\labelenumi{\alph{enumi})}
\item
  Show that every point \(x \in \mathbf{E}\) can be written uniquely in the form \(x = x_{+} - x_{-}\) where \(x_{+} \in \mathbf{V}\) and \(x_{-} \in \mathbf{V}^{\circ}\) , and \(\langle x_{+}|x_{-}\rangle = 0\) .
\item
  For every facet F of V (II, p.~87, exerc. 3), F is either the point 0 or is a convex cone with vertex 0; the set of all \(y \in V^{\circ}\) which are orthogonal to F is a closed facet \(F'\) of \(V^{\circ}\) (but this is not the « dual facet » of F in the sense of II, p.~87, exerc. 6, the latter being empty).
\item
  Take for E the set of all Hilbert-Schmidt endomorphisms of a real hilbertian space, and for V the set of positive elements in E. Show that we have \(V^{\circ} = V\) and in this case interpret the result of a) (to see that \(V \subset V^{\circ}\) , use cor. 1 of V, p.~56).
\item
  Under the hypothesis of c), let \(v \in V\) ; the set L of all \(x \in H\) such that \(\langle v(x)|x \rangle = 0\) , or, which is the same, such that \(v(x) = 0\) (V, p.~77, exerc. 10) is a closed vector subspace of H, and the facet F of v in V is the closed set of all \(u \in V\) such that \(u(x) = 0\) for all \(x \in L\) ; it can be identified with the cone of all positive Hilbert-Schmidt endomorphisms of the hilbertian space \(L^{\circ}\) . Deduce that the projection from E onto the convex set F (V, p.~11) is identical with the orthogonal projector from E onto the closed vector subspace of E generated by F.
\end{enumerate}

¶ 13) Let E be a hilbertian space, G a subgroup of the group of all automorphisms of the hilbertian space structure of E. Let \(E^{G}\) be the closed vector subspace of E consisting of all vectors invariant under G, and p the orthoprojector from E onto \(E^{G}\) .

\begin{enumerate}
\def\labelenumi{\alph{enumi})}
\item
  Show that the orthogonal complement of \(\mathbf{E}^{\mathrm{G}}\) in \(\mathbf{E}\) is the closed vector subspace generated by the vectors \(s.x - x\) , where \(s\) ranges over \(\mathbf{G}\) and \(x \in \mathbf{E}\) .
\item
  Let \(\mathbf{H}\) be a non-empty closed convex subset of \(\mathbf{E}\) which is stable under \(\mathbf{G}\) . Show that the projection of 0 onto \(\mathbf{H}\) belongs to \(\mathbf{E}^{\mathrm{G}}\) .
\item
  Suppose that H is the closed convex envelope of the orbit of a point x of E, and let a be the projection of 0 onto H. Show that \(a = p(x)\) and that \(H \cap E^{G}\) reduces to the point a (« Birkhoff-Alaoglu th. »). (Observe that x - a is contained in the orthogonal complement of \(E^{G}\) .)
\item
  Suppose G is generated by an automorphism \(u\) of E. Show that \(p(x) = \lim_{n\to \infty}\frac{1}{n + 1}\sum_{j = 0}^{n}u^{j}(x)\)
\end{enumerate}

for all \(x \in \mathrm{E}\) (if \(y_{n} = \frac{1}{n + 1} \sum_{j=0}^{n} u^{j}(x)\) , note that the sequence \((y_{n})\) has a weak limit point \(a\) , and that \(u(a) = a\) , then use \(c\) ).

\begin{enumerate}
\def\labelenumi{\alph{enumi})}
\setcounter{enumi}{4}
\tightlist
\item
  Suppose \(G\) is the image of a homomorphism \(t \mapsto u_t\) from \(\mathbf{R}\) onto the group of automorphisms of \(E\) , such that for all \(x \in E\) , \(t \mapsto u_t \cdot x\) is a continuous mapping from \(\mathbf{R}\) into \(E\) . Show that

\(p(x) = \lim_{T \to \infty} \frac{1}{T} \int_{0}^{T} u_t \cdot x \, dt\)

for all \(x \in E\) .

\end{enumerate}

\begin{enumerate}
\def\labelenumi{\alph{enumi})}
\setcounter{enumi}{5}
\tightlist
\item
  Suppose that there exists an element \(x \neq 0\) of \(\mathbf{E}\) and a number \(\alpha\) such that \(0 < \alpha < 1\) and \(\| s.x - x\| \leqslant \alpha \| x\|\) for all \(s \in \mathbf{G}\) . Show that \(\mathrm{E}^{\mathrm{G}} \neq \{0\}\) (use \(c\) ).
\end{enumerate}

\begin{enumerate}
\def\labelenumi{\arabic{enumi})}
\setcounter{enumi}{13}
\tightlist
\item
  Let E be a complex hilbertian space, and T a Hilbert-Schmidt endomorphism of E. a) Let R and L be positive Hilbert-Schmidt endomorphisms such that \(R^{2} = T^{*}T\) and \(L^{2} = TT^{*}\) (V, p.~57, cor. 3); let \(R = \text{abs}(T)\) , and call this the « absolute value » of T (cf.~V, p.~75, exerc. 4); we have \(L = \text{abs}(T^{*})\) . Show that \(\text{Ker}(T) = \text{Ker}(R)\) and \(\overline{L(\text{E})} = \overline{T(\text{E})}\) . There exists one, and only one isometry V from \(R(\text{E})\) onto \(T(\text{E})\) such that T = VR; if we extend V by continuity to \(\overline{R(\text{E})}\) , then to an operator \(U \in \mathcal{L}(\text{E})\) by taking U.x = 0 on the orthogonal complement of \(R(\text{E})\) , we also have T = UR (polar decomposition of T). Then \(R = U^{*}T = U^{*}UR = RU^{*}U\) and \(L = URU^{*}\) , \(T = LU^{*}\) . If T belongs to \(\mathcal{L}^{1}(\text{E})\) , then so does \(R = \text{abs}(T)\) , and T is the product of two Hilbert-Schmidt endomorphisms.
\end{enumerate}

\begin{enumerate}
\def\labelenumi{\alph{enumi})}
\setcounter{enumi}{1}
\item
  If \(T\) belongs to \(\mathcal{L}^1 (\mathrm{E})\) , show that \(\operatorname {Tr}(\operatorname {abs}(T)) = \sup (\sum_i |\langle a_i | T.b_i\rangle |)\) where, on the right hand side, \((a_{i})\) and \((b_{i})\) range over the set of orthonormal bases of E (use the polar decomposition of \(T\) ). Show that if we put \(\| T\| _1 = \operatorname {Tr}(\operatorname {abs}(T))\) , then \(\| T\| _1\) is a norm on the space \(\mathcal{L}^1 (\mathrm{E})\) , such that \(\| T\| _2\leqslant \| T\| _1\) .
\item
  Conversely, if \(T \in \mathcal{L}(\mathrm{E})\) is such that, for every pair \(((a_i), (b_i))\) of orthonormal bases of \(\mathbf{E}\) , the sum \(\sum_{i} |\langle a_i | T \cdot b_i \rangle|\) is finite, then \(T \in \mathcal{L}^1(\mathrm{E})\) (first observe that \(T\) is a Hilbert-Schmidt endomorphism, then use the polar decomposition of \(T\) ).
\item
  Let \((\bar{T}_{\nu})\) be a sequence of Hilbert-Schmidt endomorphisms (resp. of elements of \(\mathcal{L}^1 (\mathrm{E})\) ) such that for every pair of points \(x,y\) of E, the sequence \((\langle x|T_{\nu}.y\rangle)\) converges to \(\langle x|T.y\rangle\) , where \(T\) is a linear mapping from E into itself; in addition, assume that the sequence of norms \(\| T_{\nu}\| _2\) (resp. \(\| T_{\nu}\| _1\) ) is bounded. Show that \(T\) is a Hilbert-Schmidt endomorphism (resp. an element of \(\mathcal{L}^1 (\mathrm{E})\) ) (use \(b\) ) and \(c\) ).
\item
  Deduce from \(d\) ) that the space \(\mathcal{L}^1 (\mathrm{E})\) is a Banach space for the norm \(\| T\| _1\) .
\item
  For an endomorphism \(T \in \mathcal{L}(\mathrm{E})\) to belong to \(\mathcal{L}^1(\mathrm{E})\) it is necessary and sufficient that, for at least one orthonormal basis \((e_i)\) of \(\mathrm{E}\) , the sum \(\sum_{i} \| T.e_i\|\) is finite (with the notations of \(a\) ), note that \(|\langle e_i|R.e_i\rangle| \leqslant \| T.e_i\|\) ).

\item
  In the space \(\ell_{\mathbf{C}}^2\) , let \((e_n)\) be the canonical orthonormal basis, and let \(a = \sum_{n=0}^{\infty} \frac{1}{n+1} e_n\) ;

if \(\mathrm{F}\) is the 1-dimensional subspace \(\mathbf{C}.a\) , the orthoprojector \(p_{\mathrm{F}}\) has finite trace, but the series \(\sum_{n}\| p_{\mathrm{F}}\cdot e_n\|\) is not convergent.
\end{enumerate}

\begin{enumerate}
\def\labelenumi{\arabic{enumi})}
\setcounter{enumi}{14}
\tightlist
\item
  Let \(\mathbf{E}\) be a complex hilbertian space; let \(\mathcal{B} = \mathcal{L}(\mathbf{E})\) denote the algebra of continuous endomorphisms of \(\mathbf{E}\) , endowed with the usual norm \(\| T\| = \sup_{\| x\| \leqslant 1}\| T.x\|\) . For every pair of points \(x, y\) in E, let \(\omega_{x,y}\) denote the continuous linear form \(T \mapsto \langle x|T.y \rangle\) on \(\mathcal{B}\) , and let \(\mathcal{B}_{\circ}\) be the closed subspace of the strong dual \(\mathcal{B}'\) of the Banach space \(\mathcal{B}\) , generated by the \(\omega_{x,y}\) . a) Show that the linear mapping which associates to every \(T \in \mathcal{B}\) the linear form \(\omega \mapsto \langle \omega, T \rangle\) on \(\mathcal{B}_{\circ}\) , is an isometry from \(\mathcal{B}\) onto the strong dual of \(\mathcal{B}_{\circ}\) ; in other words, \(\mathcal{B}_{\circ}\) is a predual (IV, p.~56, exerc. 23) of \(\mathcal{B}\) .
\end{enumerate}

The topology \(\sigma (\mathcal{B},\mathcal{B}_{\circ})\) on \(\mathcal{B}\) is called the ultraweak topology.

\begin{enumerate}
\def\labelenumi{\alph{enumi})}
\setcounter{enumi}{1}
\item
  For every element \(T\) of \(\mathcal{L}^1(\mathrm{E})\) , we define a linear form \(\phi_T\) on \(\mathcal{B}\) by the formula \(\phi_T(S) = \operatorname{Tr}(ST)\) for every operator \(S \in \mathcal{B}\) . Show that \(\phi_T\) is continuous and that the mapping \(T \mapsto \phi_T\) is an isometry from the Banach space \(\mathcal{L}^1(\mathrm{E})\) (exerc. 14, e)) onto the Banach space \(\mathcal{B}_{\circ}\) (first consider the case when \(T\) has finite rank).
\item
  Let \(\mathcal{B}_{\circ \circ}\) be the vector subspace of \(\mathcal{B}_{\circ}\) generated by the \(\omega_{x,y}\) (in such a way that \(\mathcal{B}_{\circ} = \overline{\mathcal{B}}_{\circ \circ} \) ). Show that \(\mathcal{B}_{\circ \circ}\) is barrelled (note that a subset of \(\mathcal{B}\) which is bounded for \(\sigma (\mathcal{B},\mathcal{B}_{\circ \circ})\) is bounded for the norm topology).
\item
  Let \(F_{n}\) be the subspace of \(\mathcal{B}_{\circ \circ}\) which is the image of the set of all endomorphisms of E of rank \(\leqslant n\) under the isometry defined in b). Show that \(F_{n}\) is nowhere dense in \(\mathcal{B}_{\circ \circ}\) and deduce that \(\mathcal{B}_{\circ \circ}\) is not a Baire space.
\end{enumerate}

\backmatter
